\section{Introduction}
\label{sec:intro}

% fig:signflip is on page 1 (sec/0_abstract.tex); fig:overview is in the appendix.

Neural architecture search is bottlenecked by evaluation: training a candidate
costs GPU-days, and a search visits thousands. Zero-cost
proxies~\cite{abdelfattah2021zerocost} aim to reduce this bottleneck, ranking
an architecture from a single forward--backward pass at initialization, and on any
\emph{single} benchmark several rival a trained predictor~\cite{krishnakumar2022nbsz}.
They are often interpreted as \emph{capacity-associated} surrogates, tracking
parameters and FLOPs about as well as accuracy on some benchmarks
~\cite{ning2021evaluating,white2021powerful,kadlecova2024graf}, a proposed route
to cheap evaluation.

However, two observations about this zoo are in tension: many proxies correlate
with capacity, yet proxy quality is not reliably portable
~\cite{ning2021evaluating,white2021powerful,krishnakumar2022nbsz}. A proxy that
tops one benchmark or axis can be anti-correlated with accuracy on another. The
condition-number sum \px{cond-sum}, for example,
runs positive ($+0.88$) as wiring varies and negative ($-0.77$)
as width varies: the same published statistic, with opposite sign across axes
(\cref{fig:signflip}). We ask which statistical structures accompany this
transport failure and whether it changes a budgeted search outcome.

Rather than build a new proxy, we analyze the structure associated with existing
scores. The clean instrument is NAS-Bench-201~\cite{dong2021natsbench}:
its topology axis varies wiring at fixed width and its size axis varies width at
fixed wiring, so each axis turns \emph{one} knob, where a cross-\emph{benchmark}
comparison would confound family, dataset, and parameterization at once. Treating
a $29$-feature aligned zoo as one instrument, we organize each proxy's accuracy correlation
(\cref{eq:decomp}) into a \emph{capacity} term, computable label-free from the
representation spectrum, and a \emph{trainability residual} read against a measured
target $T$, the early-descent rate of the real training curve. Coincidentally,
NAS-Bench-Suite-Zero reports an ensemble adds up to $+42\%$ over the best single
proxy~\cite{krishnakumar2022nbsz}; we do not chase that number; we test how much
of a related complementarity gain remains after capacity adjustment
(\cref{app:access-evidence}). This decomposition is descriptive rather
than causal: it separates a dominant capacity-associated factor from an
early-trainability residual and then tests how both behave under transport.

Taken together, the evidence supports three bounded conclusions
(\cref{fig:results,fig:portability,fig:signflip}). First, both axes contain a
dominant capacity-associated factor, but outcome-informed orientation overstates
cross-axis loading-direction agreement. Second, a capacity-orthogonal
early-trainability association remains partly readable at initialization. Third,
several leading raw proxies reverse sign under transport, while other proxies---
including recent controls---remain positive on both axes. Concretely, we
contribute:
\begin{enumerate}\itemsep2pt
\item A \textbf{capacity--trainability decomposition} (\cref{eq:decomp}): a
dominant capacity-associated factor (\cref{sec:result1}, with the narrow
label-free yardstick of Proposition~\ref{thm:bound}) and a smaller,
capacity-orthogonal trainability association (\cref{sec:result2}) that is partly
init-accessible (\cref{sec:result3}).
\item An \textbf{empirical test} of a proposed ``label-free ceiling'' on these
benchmarks (\cref{sec:result5}): on the topology axis the best proxies largely close
the gap themselves, so the binding limitation there is portability, not magnitude;
on the size axis the residual increment sits at the noise floor.
\item A \textbf{tiered portability map} (\cref{sec:result4},
\cref{fig:portability}): three motivating within-NB201 examples, $9/13$
descriptive reversals across the broader map, and four strictly supported
cross-benchmark instances. Dimension controls narrow \px{cond-sum} to a
dimension-linked/proxy-geometry-dependent failure rather than an independent
mechanism.
\item A \textbf{controlled search consequence} (\cref{sec:searchtransfer}): two
named traps cost $+5.40$ and $+1.97$ points at $K{=}10$, while two null
directions cost zero. The target-calibrated screen is separated from target-EVAL
scoring but is not fully label-free, and we do not claim universal superiority.
\item A \textbf{reusable artifact design} (\cref{sec:artifact}): we construct a
cross-space portability map from per-architecture proxy values and
per-(proxy, space-pair) sign-stability annotations.
\end{enumerate}

\noindent Our study is \emph{observational}: we turn no controlled knobs and make
no causal claims; limitations, including the weaker size-axis arm ($n{=}138$
full-zoo), are detailed in \cref{sec:limitations}. The appendix supplies the
dependency layer under each result, organized as a per-conclusion dependency map
(\cref{app:map}).
